% !TeX root = ../教材题库_全部条目_转写工作稿.tex
% Source: XJ-B1-C05-EX-P168-Y1068
% 第5章 / 5.2 任意角的三角函数 / 教材 p.168 / PDF 第175页
% OCR未校对：公式与文字均可能有误；未完成数学TeX转写；禁止直接发给学生
% 原题及相邻原解截图（核对用）：../题目原图/XJ-B1-C05-EX-P168-Y1068-01.webp
% 本文件只在转写完成后作为题干使用；OCR 条目尚不是可靠题干。
\begin{bankitem}{5.2 例题 例题5}{第5章 / 5.2 任意角的三角函数 / 教材 p.168 / PDF 第175页}{XJ-B1-C05-EX-P168-Y1068}{OCR未校对：公式与文字均可能有误；未完成数学TeX转写；禁止直接发给学生}
\begin{ocrdraft}
[教材p.168 / PDF第175页]
\par
5
\par
例0
\par
求丷，tana的值
\par
13，
\par
解因为sin<0，sin一1，所以a是第三或第四象限角．
\par
由sin仅十cos2仅一1得
\par
COS；=1一一Sina；=「1
\par
如果是第三象限角，那么cos仅<0．于是
\par
144
\par
12
\par
COS．一
\par
169
\par
13，
\par
从而
\par
S1na
\par
5
\par
tan．一
\par
12
\par
12．
\par
如果是第四象限角，那么
\par
12
\par
5
\par
COSa．一
\par
13，
\par
tan=
\par
12．
\end{ocrdraft}
\end{bankitem}
